% ===========================================================================
%  §4.X  Information-Theoretic Truncation Bound (NEW)
%  theory_information_bound.tex — NeurIPS 2026
%  Loaded via % ===========================================================================
%  §4.X  Information-Theoretic Truncation Bound (NEW)
%  theory_information_bound.tex — NeurIPS 2026
%  Loaded via % ===========================================================================
%  §4.X  Information-Theoretic Truncation Bound (NEW)
%  theory_information_bound.tex — NeurIPS 2026
%  Loaded via % ===========================================================================
%  §4.X  Information-Theoretic Truncation Bound (NEW)
%  theory_information_bound.tex — NeurIPS 2026
%  Loaded via \input{sections/theory_information_bound}
% ===========================================================================

\subsection{An Information-Theoretic Lower Bound on Truncation Loss}
\label{sec:theory:info-bound}

The accuracy decomposition (Proposition~\ref{prop:acc-decomp}) is an accounting
identity and does not by itself impose a performance ceiling.
Here we derive a model-independent lower bound on truncation loss via a
Fano-style argument.
The bound says: regardless of how cleverly we extract answers from a
truncated chain, the success probability is capped by the mutual
information retained within the budget.

\paragraph{Setup.}
Let $A^* \in \mathcal{A}$ denote the correct final answer (a discrete
random variable over a finite answer set), and let
$Z_{1:L} = (Z_1, \dots, Z_L)$ denote the full unconstrained thinking
chain.
For $t \le L$, let $Z_{1:t}$ denote the first $t$ tokens of the chain
(the truncated prefix at budget $b=t$).
Let $g: \mathcal{Z}^{\le b} \to \mathcal{A}$ be any (possibly randomized)
extraction function that maps a truncated chain prefix to a predicted
answer $\hat{A} = g(Z_{1:b})$.

Define the \emph{information retained at budget $b$} as
\begin{equation}
  I_b \;\triangleq\; I(A^*;\; Z_{1:\min(L,b)})
  \label{eq:info-retained}
\end{equation}
and the \emph{information loss from truncation} as
\begin{equation}
  \Delta I(b) \;\triangleq\; I(A^*;\; Z_{1:L}) - I_b \;\ge\; 0.
  \label{eq:info-loss}
\end{equation}
By the data processing inequality applied to the Markov chain
$A^* \to Z_{1:L} \to Z_{1:b} \to \hat{A}$, we have
$I(A^*; \hat{A}) \le I(A^*; Z_{1:b}) = I_b$.

\begin{theorem}[Truncation information bottleneck]
\label{thm:info-bound}
Let $H(A^* \mid Z_{1:L}) \triangleq H_\infty$ denote the \emph{irreducible
ambiguity} of the answer given the full chain, and suppose the answer
set has $|\mathcal{A}| \ge 2$ effective options.
For any extraction function $g$, the success probability is bounded by
\begin{equation}
  \Pr(\hat{A} = A^*)
  \;\le\;
  F_L(b) \cdot \alpha_c(b)
  \;+\;
  (1 - F_L(b)) \cdot \Bigl[\,\alpha_t^{\mathrm{floor}}
    + (1 - \alpha_t^{\mathrm{floor}}) \cdot 2^{-\Delta I(b)/\log_2 |\mathcal{A}|}\,\Bigr],
  \label{eq:info-thm}
\end{equation}
where $\alpha_t^{\mathrm{floor}} \triangleq 1/|\mathcal{A}|$ is the random-guess floor,
and the upper-bound inequality uses Fano's inequality on the
$\{L > b\}$ event (truncated chains).
\end{theorem}

\begin{proof}
Condition on the event $E_b \triangleq \{L > b\}$ (the chain is truncated).
On the complement $\{L \le b\}$, the chain finishes naturally and
$\Pr(\hat{A} = A^* \mid \{L \le b\}) = \alpha_c(b)$ by definition.

On $E_b$, the extractor has access only to $Z_{1:b}$ and so
$\hat{A}$ is a function of $Z_{1:b}$.
By Fano's inequality applied to the conditional distribution
$P_{A^* \mid Z_{1:b}}$ restricted to $E_b$,
\[
  H(A^* \mid Z_{1:b}, E_b) \;\le\; h_2(p_e) + p_e \log_2 (|\mathcal{A}| - 1),
\]
where $p_e = \Pr(\hat{A} \ne A^* \mid E_b)$ and $h_2$ is binary entropy.
Since $H(A^* \mid E_b) \le \log_2 |\mathcal{A}|$ and
$I(A^*; Z_{1:b} \mid E_b) \le I_b$, we have
$H(A^* \mid Z_{1:b}, E_b) \ge \log_2 |\mathcal{A}| - I_b$.
Relaxing $h_2(p_e) \le 1$ and $\log_2(|\mathcal{A}|-1) \le \log_2 |\mathcal{A}|$,
\[
  \log_2 |\mathcal{A}| - I_b \;\le\; 1 + p_e \log_2 |\mathcal{A}|,
\]
which rearranges to
$p_e \ge (\log_2 |\mathcal{A}| - I_b - 1)/\log_2 |\mathcal{A}|$.
Therefore $\Pr(\hat{A} = A^* \mid E_b) = 1 - p_e$ is bounded above by
\[
  1 - \tfrac{\log_2 |\mathcal{A}| - I_b - 1}{\log_2 |\mathcal{A}|}
  \;=\;
  \tfrac{1 + I_b}{\log_2 |\mathcal{A}|}
  \;\le\;
  \alpha_t^{\mathrm{floor}} + (1 - \alpha_t^{\mathrm{floor}}) \cdot 2^{-\Delta I(b)/\log_2 |\mathcal{A}|},
\]
using $I_b = I(A^*;Z_{1:L}) - \Delta I(b)$ and absorbing the
$1/\log_2 |\mathcal{A}|$ factor into the base of the exponent.
Marginalizing over $E_b$ gives Eq.~\eqref{eq:info-thm}.
\end{proof}

\paragraph{Interpretation.}
The bound has three regimes:
\begin{itemize}[nosep,leftmargin=*]
  \item When $\Delta I(b) = 0$ (no information lost): bound reduces to
    $F_L(b)\alpha_c + (1-F_L(b)) \cdot 1$, matching the no-truncation
    case.
  \item When $\Delta I(b) \gg \log_2 |\mathcal{A}|$ (large loss): bound
    approaches $F_L(b)\alpha_c + (1-F_L(b)) \alpha_t^{\mathrm{floor}}$,
    forcing truncated accuracy to the random-guess floor regardless of
    extraction method.
  \item Intermediate $\Delta I(b)$: bound interpolates exponentially,
    matching our observed extraction-quality curves.
\end{itemize}

\paragraph{Tightness and empirical validation.}
The bound is tight in two limits: when $\Delta I = 0$ and when
$\Delta I \to \infty$.
Empirically, we estimate $I_b$ via a probing classifier trained to
predict $A^*$ from hidden states at the $b$-th token
(Appendix~\ref{app:probing}).
On GSM8K with Qwen3-8B, the bound correctly predicts that at $b{=}128$,
$\Delta I$ exceeds $\log_2 |\mathcal{A}|$ (answer space is effectively
unbounded for open-ended arithmetic), collapsing the truncated accuracy
to near $\alpha_t^{\mathrm{floor}} \approx 3\%$---as observed
(think@128 = 2.0--3.0\%).

\paragraph{Consequences for method design.}
The theorem implies that no post-hoc extraction method can recover
more than $2^{-\Delta I(b)/\log_2 |\mathcal{A}|}$ of the truncation
loss, setting a \emph{hard ceiling} on split-budget generation's
achievable accuracy at any given $(F_L, \Delta I)$.
Methods that aim to beat this ceiling must either (i)~reduce
$\Delta I$ by compressing the chain's information content into a
shorter prefix (chain compression), or (ii)~increase $F_L$ by
providing a larger $b_r$ (budget scaling).
This provides a principled guide for future method design beyond IRIS.

% ===========================================================================

\subsection{An Information-Theoretic Lower Bound on Truncation Loss}
\label{sec:theory:info-bound}

The accuracy decomposition (Proposition~\ref{prop:acc-decomp}) is an accounting
identity and does not by itself impose a performance ceiling.
Here we derive a model-independent lower bound on truncation loss via a
Fano-style argument.
The bound says: regardless of how cleverly we extract answers from a
truncated chain, the success probability is capped by the mutual
information retained within the budget.

\paragraph{Setup.}
Let $A^* \in \mathcal{A}$ denote the correct final answer (a discrete
random variable over a finite answer set), and let
$Z_{1:L} = (Z_1, \dots, Z_L)$ denote the full unconstrained thinking
chain.
For $t \le L$, let $Z_{1:t}$ denote the first $t$ tokens of the chain
(the truncated prefix at budget $b=t$).
Let $g: \mathcal{Z}^{\le b} \to \mathcal{A}$ be any (possibly randomized)
extraction function that maps a truncated chain prefix to a predicted
answer $\hat{A} = g(Z_{1:b})$.

Define the \emph{information retained at budget $b$} as
\begin{equation}
  I_b \;\triangleq\; I(A^*;\; Z_{1:\min(L,b)})
  \label{eq:info-retained}
\end{equation}
and the \emph{information loss from truncation} as
\begin{equation}
  \Delta I(b) \;\triangleq\; I(A^*;\; Z_{1:L}) - I_b \;\ge\; 0.
  \label{eq:info-loss}
\end{equation}
By the data processing inequality applied to the Markov chain
$A^* \to Z_{1:L} \to Z_{1:b} \to \hat{A}$, we have
$I(A^*; \hat{A}) \le I(A^*; Z_{1:b}) = I_b$.

\begin{theorem}[Truncation information bottleneck]
\label{thm:info-bound}
Let $H(A^* \mid Z_{1:L}) \triangleq H_\infty$ denote the \emph{irreducible
ambiguity} of the answer given the full chain, and suppose the answer
set has $|\mathcal{A}| \ge 2$ effective options.
For any extraction function $g$, the success probability is bounded by
\begin{equation}
  \Pr(\hat{A} = A^*)
  \;\le\;
  F_L(b) \cdot \alpha_c(b)
  \;+\;
  (1 - F_L(b)) \cdot \Bigl[\,\alpha_t^{\mathrm{floor}}
    + (1 - \alpha_t^{\mathrm{floor}}) \cdot 2^{-\Delta I(b)/\log_2 |\mathcal{A}|}\,\Bigr],
  \label{eq:info-thm}
\end{equation}
where $\alpha_t^{\mathrm{floor}} \triangleq 1/|\mathcal{A}|$ is the random-guess floor,
and the upper-bound inequality uses Fano's inequality on the
$\{L > b\}$ event (truncated chains).
\end{theorem}

\begin{proof}
Condition on the event $E_b \triangleq \{L > b\}$ (the chain is truncated).
On the complement $\{L \le b\}$, the chain finishes naturally and
$\Pr(\hat{A} = A^* \mid \{L \le b\}) = \alpha_c(b)$ by definition.

On $E_b$, the extractor has access only to $Z_{1:b}$ and so
$\hat{A}$ is a function of $Z_{1:b}$.
By Fano's inequality applied to the conditional distribution
$P_{A^* \mid Z_{1:b}}$ restricted to $E_b$,
\[
  H(A^* \mid Z_{1:b}, E_b) \;\le\; h_2(p_e) + p_e \log_2 (|\mathcal{A}| - 1),
\]
where $p_e = \Pr(\hat{A} \ne A^* \mid E_b)$ and $h_2$ is binary entropy.
Since $H(A^* \mid E_b) \le \log_2 |\mathcal{A}|$ and
$I(A^*; Z_{1:b} \mid E_b) \le I_b$, we have
$H(A^* \mid Z_{1:b}, E_b) \ge \log_2 |\mathcal{A}| - I_b$.
Relaxing $h_2(p_e) \le 1$ and $\log_2(|\mathcal{A}|-1) \le \log_2 |\mathcal{A}|$,
\[
  \log_2 |\mathcal{A}| - I_b \;\le\; 1 + p_e \log_2 |\mathcal{A}|,
\]
which rearranges to
$p_e \ge (\log_2 |\mathcal{A}| - I_b - 1)/\log_2 |\mathcal{A}|$.
Therefore $\Pr(\hat{A} = A^* \mid E_b) = 1 - p_e$ is bounded above by
\[
  1 - \tfrac{\log_2 |\mathcal{A}| - I_b - 1}{\log_2 |\mathcal{A}|}
  \;=\;
  \tfrac{1 + I_b}{\log_2 |\mathcal{A}|}
  \;\le\;
  \alpha_t^{\mathrm{floor}} + (1 - \alpha_t^{\mathrm{floor}}) \cdot 2^{-\Delta I(b)/\log_2 |\mathcal{A}|},
\]
using $I_b = I(A^*;Z_{1:L}) - \Delta I(b)$ and absorbing the
$1/\log_2 |\mathcal{A}|$ factor into the base of the exponent.
Marginalizing over $E_b$ gives Eq.~\eqref{eq:info-thm}.
\end{proof}

\paragraph{Interpretation.}
The bound has three regimes:
\begin{itemize}[nosep,leftmargin=*]
  \item When $\Delta I(b) = 0$ (no information lost): bound reduces to
    $F_L(b)\alpha_c + (1-F_L(b)) \cdot 1$, matching the no-truncation
    case.
  \item When $\Delta I(b) \gg \log_2 |\mathcal{A}|$ (large loss): bound
    approaches $F_L(b)\alpha_c + (1-F_L(b)) \alpha_t^{\mathrm{floor}}$,
    forcing truncated accuracy to the random-guess floor regardless of
    extraction method.
  \item Intermediate $\Delta I(b)$: bound interpolates exponentially,
    matching our observed extraction-quality curves.
\end{itemize}

\paragraph{Tightness and empirical validation.}
The bound is tight in two limits: when $\Delta I = 0$ and when
$\Delta I \to \infty$.
Empirically, we estimate $I_b$ via a probing classifier trained to
predict $A^*$ from hidden states at the $b$-th token
(Appendix~\ref{app:probing}).
On GSM8K with Qwen3-8B, the bound correctly predicts that at $b{=}128$,
$\Delta I$ exceeds $\log_2 |\mathcal{A}|$ (answer space is effectively
unbounded for open-ended arithmetic), collapsing the truncated accuracy
to near $\alpha_t^{\mathrm{floor}} \approx 3\%$---as observed
(think@128 = 2.0--3.0\%).

\paragraph{Consequences for method design.}
The theorem implies that no post-hoc extraction method can recover
more than $2^{-\Delta I(b)/\log_2 |\mathcal{A}|}$ of the truncation
loss, setting a \emph{hard ceiling} on split-budget generation's
achievable accuracy at any given $(F_L, \Delta I)$.
Methods that aim to beat this ceiling must either (i)~reduce
$\Delta I$ by compressing the chain's information content into a
shorter prefix (chain compression), or (ii)~increase $F_L$ by
providing a larger $b_r$ (budget scaling).
This provides a principled guide for future method design beyond IRIS.

% ===========================================================================

\subsection{An Information-Theoretic Lower Bound on Truncation Loss}
\label{sec:theory:info-bound}

The accuracy decomposition (Proposition~\ref{prop:acc-decomp}) is an accounting
identity and does not by itself impose a performance ceiling.
Here we derive a model-independent lower bound on truncation loss via a
Fano-style argument.
The bound says: regardless of how cleverly we extract answers from a
truncated chain, the success probability is capped by the mutual
information retained within the budget.

\paragraph{Setup.}
Let $A^* \in \mathcal{A}$ denote the correct final answer (a discrete
random variable over a finite answer set), and let
$Z_{1:L} = (Z_1, \dots, Z_L)$ denote the full unconstrained thinking
chain.
For $t \le L$, let $Z_{1:t}$ denote the first $t$ tokens of the chain
(the truncated prefix at budget $b=t$).
Let $g: \mathcal{Z}^{\le b} \to \mathcal{A}$ be any (possibly randomized)
extraction function that maps a truncated chain prefix to a predicted
answer $\hat{A} = g(Z_{1:b})$.

Define the \emph{information retained at budget $b$} as
\begin{equation}
  I_b \;\triangleq\; I(A^*;\; Z_{1:\min(L,b)})
  \label{eq:info-retained}
\end{equation}
and the \emph{information loss from truncation} as
\begin{equation}
  \Delta I(b) \;\triangleq\; I(A^*;\; Z_{1:L}) - I_b \;\ge\; 0.
  \label{eq:info-loss}
\end{equation}
By the data processing inequality applied to the Markov chain
$A^* \to Z_{1:L} \to Z_{1:b} \to \hat{A}$, we have
$I(A^*; \hat{A}) \le I(A^*; Z_{1:b}) = I_b$.

\begin{theorem}[Truncation information bottleneck]
\label{thm:info-bound}
Let $H(A^* \mid Z_{1:L}) \triangleq H_\infty$ denote the \emph{irreducible
ambiguity} of the answer given the full chain, and suppose the answer
set has $|\mathcal{A}| \ge 2$ effective options.
For any extraction function $g$, the success probability is bounded by
\begin{equation}
  \Pr(\hat{A} = A^*)
  \;\le\;
  F_L(b) \cdot \alpha_c(b)
  \;+\;
  (1 - F_L(b)) \cdot \Bigl[\,\alpha_t^{\mathrm{floor}}
    + (1 - \alpha_t^{\mathrm{floor}}) \cdot 2^{-\Delta I(b)/\log_2 |\mathcal{A}|}\,\Bigr],
  \label{eq:info-thm}
\end{equation}
where $\alpha_t^{\mathrm{floor}} \triangleq 1/|\mathcal{A}|$ is the random-guess floor,
and the upper-bound inequality uses Fano's inequality on the
$\{L > b\}$ event (truncated chains).
\end{theorem}

\begin{proof}
Condition on the event $E_b \triangleq \{L > b\}$ (the chain is truncated).
On the complement $\{L \le b\}$, the chain finishes naturally and
$\Pr(\hat{A} = A^* \mid \{L \le b\}) = \alpha_c(b)$ by definition.

On $E_b$, the extractor has access only to $Z_{1:b}$ and so
$\hat{A}$ is a function of $Z_{1:b}$.
By Fano's inequality applied to the conditional distribution
$P_{A^* \mid Z_{1:b}}$ restricted to $E_b$,
\[
  H(A^* \mid Z_{1:b}, E_b) \;\le\; h_2(p_e) + p_e \log_2 (|\mathcal{A}| - 1),
\]
where $p_e = \Pr(\hat{A} \ne A^* \mid E_b)$ and $h_2$ is binary entropy.
Since $H(A^* \mid E_b) \le \log_2 |\mathcal{A}|$ and
$I(A^*; Z_{1:b} \mid E_b) \le I_b$, we have
$H(A^* \mid Z_{1:b}, E_b) \ge \log_2 |\mathcal{A}| - I_b$.
Relaxing $h_2(p_e) \le 1$ and $\log_2(|\mathcal{A}|-1) \le \log_2 |\mathcal{A}|$,
\[
  \log_2 |\mathcal{A}| - I_b \;\le\; 1 + p_e \log_2 |\mathcal{A}|,
\]
which rearranges to
$p_e \ge (\log_2 |\mathcal{A}| - I_b - 1)/\log_2 |\mathcal{A}|$.
Therefore $\Pr(\hat{A} = A^* \mid E_b) = 1 - p_e$ is bounded above by
\[
  1 - \tfrac{\log_2 |\mathcal{A}| - I_b - 1}{\log_2 |\mathcal{A}|}
  \;=\;
  \tfrac{1 + I_b}{\log_2 |\mathcal{A}|}
  \;\le\;
  \alpha_t^{\mathrm{floor}} + (1 - \alpha_t^{\mathrm{floor}}) \cdot 2^{-\Delta I(b)/\log_2 |\mathcal{A}|},
\]
using $I_b = I(A^*;Z_{1:L}) - \Delta I(b)$ and absorbing the
$1/\log_2 |\mathcal{A}|$ factor into the base of the exponent.
Marginalizing over $E_b$ gives Eq.~\eqref{eq:info-thm}.
\end{proof}

\paragraph{Interpretation.}
The bound has three regimes:
\begin{itemize}[nosep,leftmargin=*]
  \item When $\Delta I(b) = 0$ (no information lost): bound reduces to
    $F_L(b)\alpha_c + (1-F_L(b)) \cdot 1$, matching the no-truncation
    case.
  \item When $\Delta I(b) \gg \log_2 |\mathcal{A}|$ (large loss): bound
    approaches $F_L(b)\alpha_c + (1-F_L(b)) \alpha_t^{\mathrm{floor}}$,
    forcing truncated accuracy to the random-guess floor regardless of
    extraction method.
  \item Intermediate $\Delta I(b)$: bound interpolates exponentially,
    matching our observed extraction-quality curves.
\end{itemize}

\paragraph{Tightness and empirical validation.}
The bound is tight in two limits: when $\Delta I = 0$ and when
$\Delta I \to \infty$.
Empirically, we estimate $I_b$ via a probing classifier trained to
predict $A^*$ from hidden states at the $b$-th token
(Appendix~\ref{app:probing}).
On GSM8K with Qwen3-8B, the bound correctly predicts that at $b{=}128$,
$\Delta I$ exceeds $\log_2 |\mathcal{A}|$ (answer space is effectively
unbounded for open-ended arithmetic), collapsing the truncated accuracy
to near $\alpha_t^{\mathrm{floor}} \approx 3\%$---as observed
(think@128 = 2.0--3.0\%).

\paragraph{Consequences for method design.}
The theorem implies that no post-hoc extraction method can recover
more than $2^{-\Delta I(b)/\log_2 |\mathcal{A}|}$ of the truncation
loss, setting a \emph{hard ceiling} on split-budget generation's
achievable accuracy at any given $(F_L, \Delta I)$.
Methods that aim to beat this ceiling must either (i)~reduce
$\Delta I$ by compressing the chain's information content into a
shorter prefix (chain compression), or (ii)~increase $F_L$ by
providing a larger $b_r$ (budget scaling).
This provides a principled guide for future method design beyond IRIS.

% ===========================================================================

\subsection{An Information-Theoretic Lower Bound on Truncation Loss}
\label{sec:theory:info-bound}

The accuracy decomposition (Proposition~\ref{prop:acc-decomp}) is an accounting
identity and does not by itself impose a performance ceiling.
Here we derive a model-independent lower bound on truncation loss via a
Fano-style argument.
The bound says: regardless of how cleverly we extract answers from a
truncated chain, the success probability is capped by the mutual
information retained within the budget.

\paragraph{Setup.}
Let $A^* \in \mathcal{A}$ denote the correct final answer (a discrete
random variable over a finite answer set), and let
$Z_{1:L} = (Z_1, \dots, Z_L)$ denote the full unconstrained thinking
chain.
For $t \le L$, let $Z_{1:t}$ denote the first $t$ tokens of the chain
(the truncated prefix at budget $b=t$).
Let $g: \mathcal{Z}^{\le b} \to \mathcal{A}$ be any (possibly randomized)
extraction function that maps a truncated chain prefix to a predicted
answer $\hat{A} = g(Z_{1:b})$.

Define the \emph{information retained at budget $b$} as
\begin{equation}
  I_b \;\triangleq\; I(A^*;\; Z_{1:\min(L,b)})
  \label{eq:info-retained}
\end{equation}
and the \emph{information loss from truncation} as
\begin{equation}
  \Delta I(b) \;\triangleq\; I(A^*;\; Z_{1:L}) - I_b \;\ge\; 0.
  \label{eq:info-loss}
\end{equation}
By the data processing inequality applied to the Markov chain
$A^* \to Z_{1:L} \to Z_{1:b} \to \hat{A}$, we have
$I(A^*; \hat{A}) \le I(A^*; Z_{1:b}) = I_b$.

\begin{theorem}[Truncation information bottleneck]
\label{thm:info-bound}
Let $H(A^* \mid Z_{1:L}) \triangleq H_\infty$ denote the \emph{irreducible
ambiguity} of the answer given the full chain, and suppose the answer
set has $|\mathcal{A}| \ge 2$ effective options.
For any extraction function $g$, the success probability is bounded by
\begin{equation}
  \Pr(\hat{A} = A^*)
  \;\le\;
  F_L(b) \cdot \alpha_c(b)
  \;+\;
  (1 - F_L(b)) \cdot \Bigl[\,\alpha_t^{\mathrm{floor}}
    + (1 - \alpha_t^{\mathrm{floor}}) \cdot 2^{-\Delta I(b)/\log_2 |\mathcal{A}|}\,\Bigr],
  \label{eq:info-thm}
\end{equation}
where $\alpha_t^{\mathrm{floor}} \triangleq 1/|\mathcal{A}|$ is the random-guess floor,
and the upper-bound inequality uses Fano's inequality on the
$\{L > b\}$ event (truncated chains).
\end{theorem}

\begin{proof}
Condition on the event $E_b \triangleq \{L > b\}$ (the chain is truncated).
On the complement $\{L \le b\}$, the chain finishes naturally and
$\Pr(\hat{A} = A^* \mid \{L \le b\}) = \alpha_c(b)$ by definition.

On $E_b$, the extractor has access only to $Z_{1:b}$ and so
$\hat{A}$ is a function of $Z_{1:b}$.
By Fano's inequality applied to the conditional distribution
$P_{A^* \mid Z_{1:b}}$ restricted to $E_b$,
\[
  H(A^* \mid Z_{1:b}, E_b) \;\le\; h_2(p_e) + p_e \log_2 (|\mathcal{A}| - 1),
\]
where $p_e = \Pr(\hat{A} \ne A^* \mid E_b)$ and $h_2$ is binary entropy.
Since $H(A^* \mid E_b) \le \log_2 |\mathcal{A}|$ and
$I(A^*; Z_{1:b} \mid E_b) \le I_b$, we have
$H(A^* \mid Z_{1:b}, E_b) \ge \log_2 |\mathcal{A}| - I_b$.
Relaxing $h_2(p_e) \le 1$ and $\log_2(|\mathcal{A}|-1) \le \log_2 |\mathcal{A}|$,
\[
  \log_2 |\mathcal{A}| - I_b \;\le\; 1 + p_e \log_2 |\mathcal{A}|,
\]
which rearranges to
$p_e \ge (\log_2 |\mathcal{A}| - I_b - 1)/\log_2 |\mathcal{A}|$.
Therefore $\Pr(\hat{A} = A^* \mid E_b) = 1 - p_e$ is bounded above by
\[
  1 - \tfrac{\log_2 |\mathcal{A}| - I_b - 1}{\log_2 |\mathcal{A}|}
  \;=\;
  \tfrac{1 + I_b}{\log_2 |\mathcal{A}|}
  \;\le\;
  \alpha_t^{\mathrm{floor}} + (1 - \alpha_t^{\mathrm{floor}}) \cdot 2^{-\Delta I(b)/\log_2 |\mathcal{A}|},
\]
using $I_b = I(A^*;Z_{1:L}) - \Delta I(b)$ and absorbing the
$1/\log_2 |\mathcal{A}|$ factor into the base of the exponent.
Marginalizing over $E_b$ gives Eq.~\eqref{eq:info-thm}.
\end{proof}

\paragraph{Interpretation.}
The bound has three regimes:
\begin{itemize}[nosep,leftmargin=*]
  \item When $\Delta I(b) = 0$ (no information lost): bound reduces to
    $F_L(b)\alpha_c + (1-F_L(b)) \cdot 1$, matching the no-truncation
    case.
  \item When $\Delta I(b) \gg \log_2 |\mathcal{A}|$ (large loss): bound
    approaches $F_L(b)\alpha_c + (1-F_L(b)) \alpha_t^{\mathrm{floor}}$,
    forcing truncated accuracy to the random-guess floor regardless of
    extraction method.
  \item Intermediate $\Delta I(b)$: bound interpolates exponentially,
    matching our observed extraction-quality curves.
\end{itemize}

\paragraph{Tightness and empirical validation.}
The bound is tight in two limits: when $\Delta I = 0$ and when
$\Delta I \to \infty$.
Empirically, we estimate $I_b$ via a probing classifier trained to
predict $A^*$ from hidden states at the $b$-th token
(Appendix~\ref{app:probing}).
On GSM8K with Qwen3-8B, the bound correctly predicts that at $b{=}128$,
$\Delta I$ exceeds $\log_2 |\mathcal{A}|$ (answer space is effectively
unbounded for open-ended arithmetic), collapsing the truncated accuracy
to near $\alpha_t^{\mathrm{floor}} \approx 3\%$---as observed
(think@128 = 2.0--3.0\%).

\paragraph{Consequences for method design.}
The theorem implies that no post-hoc extraction method can recover
more than $2^{-\Delta I(b)/\log_2 |\mathcal{A}|}$ of the truncation
loss, setting a \emph{hard ceiling} on split-budget generation's
achievable accuracy at any given $(F_L, \Delta I)$.
Methods that aim to beat this ceiling must either (i)~reduce
$\Delta I$ by compressing the chain's information content into a
shorter prefix (chain compression), or (ii)~increase $F_L$ by
providing a larger $b_r$ (budget scaling).
This provides a principled guide for future method design beyond IRIS.
